\begin{abstract}
Current TSAD benchmarks ask whether a score predicts anomalies, but not what it adds beyond simple statistical signals. We introduce conditional detection utility (CDU), the held-out log-loss reduction from adding a detector score to a statistical reference. We evaluate fixed anomaly-score outputs from nine detectors on 350 univariate series spanning 23 source groups, using leave-one-source-out evaluation. Under the same spline probe, loss, and source weighting, MOMENT-FT and MOMENT-ZS rank fourth and fifth by detector-only utility but seventh and ninth by CDU; TranAD rises from sixth to fourth. POLY and MOMENT-ZS have nearly identical benchmark VUS-PR (0.389 versus 0.383), yet differ by 0.00593 bits of CDU. MOMENT-ZS scores are highly reconstructible from the reference ($R^2=0.904$), but rank ninth by CDU; across all nine detectors, reconstructibility and CDU have little rank correspondence (Spearman $\rho=-0.167$). Three probe families preserve the top-two order and top-four membership; five sampling seeds give the same spline ranking, and reference variants preserve the leading group. Standalone performance and beyond-reference contribution are empirically non-interchangeable evaluation axes; CDU makes the latter measurable for heterogeneous TSAD scores.
\end{abstract}
